\subsection{The politico-economic equilibrium with LOG preferences}\label{\refsec:PEELOG}
Let $\upsilon_{1,t}^i$ and $\upsilon_{2,t}^i$ denote the indirect utilities of different households (types and generations). Policy is set without commitment, with elections each $t$ aggregating preferences using probabilistic voting that maximizes a weighted sum of indirect utilities. Let $\tau^{t+1}(z_t)$ denote a continuation policy with $z_t$ being a vector of states. The political problem can then be written as
\begin{align*}
    &\max_{\tau_t}\mbox{ }\mathcal{W}_t(z_{t-1}) \equiv \omega\sum_i\left(\gamma_{t-1,i} p_{t-1} \mu_{t-1,i} v_{2,t}^i\right)+\nu_t\sum_i\left(\gamma_{t,i}\mu_{t,i} v_{1,t}^i\right) \\
    &\text{s.t. competitive equilibrium and }\tau_{t+1} = \tau^{t+1}(z_t),
\end{align*}
where $\tau^{t+1}(z_t)$ is identified by this process at $t+1$, $\omega$ is a weight reflecting relative political power of the older generation, and $\mu_{t,i}$ reflects relative voting shares for different income groups (the old generation at $t$ votes with the voting shares $\mu_{t-1,i}$ they had when young, i.e.\ shares are attached to a generation, not a period). The voting shares matter for the US in particular, where participation rises steeply with income.

With $\rho = 1$, the discount function simplifies to $B_{t+1}^i = \beta_{t,i}$ and the politico-economic equilibrium can be characterized in closed form without any continuation policy; we can exploit that all consumption functions are log-separable in $s_{t-1}$ to show that there are no endogenous states, i.e.\ $z_t = \emptyset$ (see the paper for an elaboration of this claim). In general, the first order condition for $\tau_t$ is given by:
\begin{align*}
    \omega\sum_{i}\left(\gamma_{t-1,i}p_{t-1}\mu_{t-1,i} \frac{d\upsilon_{2,t}^i}{d\tau_t}\right)+
    \nu_t\sum_i\left(\gamma_{t,i}\mu_{t,i} \frac{d\upsilon_{1,t}^i}{d\tau_t}\right) = 0,
\end{align*}
where the total derivatives (with ``hard d'') include the partial effects of changing policy directly (through economic equilibrium functions), but also the indirect effects working through the continuation policy.

\begin{subequations}\label{\refeq:PEELOG}
For the log case, where there are no effects through the continuation policy, this simplifies significantly. Using the Euler condition we have $\upsilon_{1,t}^i = (1+\beta_{t,i})\ln\left(\tilde{c}_{1,t}^i\right)+\beta_{t,i} \ln\left(R_{t+1}\right)+\beta_{t,i}\ln(\beta_{t,i}/p_t)$. Using that $\tilde{c}_{1,t}^i$ is log-separable in $s_{t-1}$ and the formula for the interest rate, it is straightforward to verify that this reduces to
\begin{align}\label{\refeq:PEELOG:v1i}
    \frac{d\upsilon_{1,t}^i}{d\tau_t} = -\frac{1}{1-\tau_t}\frac{1+\xi}{1+\alpha\xi}\left(1+\beta_{t,i}\frac{\alpha(1+\xi)}{1+\alpha\xi}\right).
\end{align}
The two channels pull in opposite directions and the second is the weaker: a higher $\tau_t$ lowers current net income, but by depressing savings it also raises the return on what is saved. For the older cohorts, the effect works through $c_{2,t}^i$ and, again being log-separable in $s_{t-1}$, this simplifies to
\begin{align}\label{\refeq:PEELOG:v2i}
    \frac{d\upsilon_{2,t}^i}{d\tau_t} =& (1-\alpha)\frac{\partial \ln\left(\Theta_{h,t}\right)}{\partial \tau_t}+\frac{\frac{1-\alpha}{\alpha}\left(1+\theta_t\left[\frac{\eta_{t-1,i}^{1+\xi}}{X_{t-1,i}^{\xi}}\frac{1}{\Gamma_{h,t-1}}-1\right]\right)}{\frac{s_{t-1,i}}{s_{t-1}}+\frac{1-\alpha}{\alpha}\tau_t\left(1+\theta_t\left[\frac{\eta_{t-1,i}^{1+\xi}}{X_{t-1,i}^{\xi}}\frac{1}{\Gamma_{h,t-1}}-1\right]\right)},
\end{align}
where
\begin{align}\label{\refeq:PEELOG:dThetah}
    \frac{\partial \ln\left(\Theta_{h,t}\right)}{\partial \tau_t} = -\frac{1}{1-\tau_t}\frac{\xi}{1+\alpha\xi}.
\end{align}
The bracket $\eta_{t-1,i}^{1+\xi}/(X_{t-1,i}^{\xi}\Gamma_{h,t-1})-1$ is the type's relative contribution base net of the average, and it is what makes the retired bloc disagree internally: with $\theta_t<1$ the numerator is larger for low-income retirees, so a Beveridgean system turns the old into stronger supporters of taxation the poorer they are.
\end{subequations}


\smalltitle{The first order condition decouples across time.}
Inspecting \eqref{\refeq:PEELOG}, every term is a function of $\tau_t$ alone. The young term \eqref{\refeq:PEELOG:v1i} contains no other policy. In the old term \eqref{\refeq:PEELOG:v2i}, the labour-supply channel enters as the \emph{partial} derivative \eqref{\refeq:PEELOG:dThetah}, which is closed form and free of $\tau_{t+1}$, while the remaining state $s_{t-1,i}/s_{t-1}$ is by \eqref{\refeq:EE:si_s}, shifted one period back, a closed-form function of $\tau_t$, $\theta_t$ and $\bm{B}_t = \bm{\beta}_{t-1}$. Writing $z_t$ for the first order condition, we therefore have
\begin{align}\label{\refeq:PEELOG:decoupling}
    z_t = z_t\left(\tau_t\right),
\end{align}
so the equilibrium path is not a recursion at all but $T$ independent scalar problems, one per period, which may be solved in any order or in parallel. This is stronger than what holds in the models with an informal sector: there the indirect utility of the old informal household depends on the level of $\Theta_{h,t}$ rather than only on its derivative, and $\Theta_{h,t}$ depends on $\tau_{t+1}$ through $\Gamma_{s,t}$, leaving a triangular system $z_t(\tau_t,\tau_{t+1})$ that must be solved backwards. Time still enters \eqref{\refeq:PEELOG:decoupling} through the vintages $\nu_t, \gamma_{t,i}, \mu_{t,i}, p_t, \theta_t$ --- demographics remain a driver of the equilibrium tax rate --- but no lag or lead of $\tau$ does. The property relies on $\theta_t$ being exogenous; if pension characteristics are chosen alongside $\tau_t$, it has to be revisited.

\smalltitle{Policy inherits both invariances.}
Under the rescaling of \eqref{\refeq:scaleInvariance}, both terms of \eqref{\refeq:PEELOG:v2i} are unchanged: the ratio $\eta_{t-1,i}^{1+\xi}/(X_{t-1,i}^{\xi}\Gamma_{h,t-1}) = y_{t-1,i}^{\eta}/\Gamma_{h,t-1}$ is a quotient of two objects that scale alike, and $s_{t-1,i}/s_{t-1}$ is invariant. Equation \eqref{\refeq:PEELOG:v1i} contains no parameters that scale at all. Hence $z_t$, and with it the equilibrium $\tau_t$, is invariant to $\lambda$. The same holds under CRRA preferences below, where every consumption level entering the first order condition scales by $\lambda$ and the common factor $\lambda^{1-1/\rho}$ divides out of $z_t = 0$.

Invariance to the hours unit $\mu$ of \eqref{\refeq:hoursUnit} is more immediate: the first order condition involves $\eta$ and $X$ only through $y_{t-1,i}^{\eta}/\Gamma_{h,t-1}$, and $y^{\eta}$ is by construction what \eqref{\refeq:hoursUnit} holds fixed. Equilibrium policy is thus a function of the shape of the productivity distribution alone --- neither its scale nor the units in which hours are counted.



\subsection{The politico-economic equilibrium with CRRA preferences}\label{\refsec:PEE}
With CRRA preferences, the general first order condition is the same, but the total derivatives of indirect utilities change significantly: the past savings level $s_{t-1}$ becomes a state variable, the continuation policy re-enters, and consumption levels acquire a role as marginal-utility weights. As we show in \cref{\refnumsec}, we will rely on numerical methods that approximate these derivatives in combination with grid searches. For the young generations, this means that we only have to be able to compute indirect utilities. For the older generations, however, we need to go one step deeper, in order to ensure that we do not need to evaluate on a large grid of the past distribution of savings $s_{t-1,i}/s_{t-1}$ (see the main paper for a thorough explanation).

\begin{subequations}\label{\refeq:PEE}
Starting from the Euler condition for the young households, we have that
\begin{align}
    \upsilon_{1,t}^i = (1+B_{t+1}^i)\frac{\left(\tilde{c}_{1,t}^i\right)^{1-1/\rho}}{1-1/\rho}.
\end{align}

Because the policy maker has to take past choices as given (such as savings and relative savings rates), we have to rely on first order conditions to efficiently identify their choices (to avoid searching over a full grid of states $s_{t-1,i}/s_{t-1}$). Whereas we can numerically approximate derivatives of $\upsilon_{1,t}^i$ using simple grids, here we have to be a bit more careful. In particular, we have to use:
\begin{align}
    \dfrac{d v_{2,t}^i}{d \tau_t} =& \left(c_{2,t}^i\right)^{1-1/\rho}\dfrac{d\ln\left(c_{2,t}^i\right)}{d\tau_t} \\
    \frac{d\ln\left(c_{2,t}^i\right)}{d\tau_t} =& (1-\alpha)\frac{d \ln\left(h_t\right)}{d \tau_t}+\frac{\frac{1-\alpha}{\alpha}\left(1+\theta_t\left[\frac{\eta_{t-1,i}^{1+\xi}}{X_{t-1,i}^{\xi}}\frac{1}{\Gamma_{h,t-1}}-1\right]\right)}{\frac{s_{t-1,i}}{s_{t-1}}+\frac{1-\alpha}{\alpha}\tau_t\left(1+\theta_t\left[\frac{\eta_{t-1,i}^{1+\xi}}{X_{t-1,i}^{\xi}}\frac{1}{\Gamma_{h,t-1}}-1\right]\right)},
\end{align}
where the second line explicitly uses that $s_{t-1,i}/s_{t-1}$ is predetermined. However, as this ratio has a closed-form expression that relies on $\tau_t$, we do not have to evaluate on the entire grid, rather only on the selection that are actually consistent with economic equilibrium conditions. This is the exact trick that allows us to only use $s_{t-1}$ as a relevant state and not the entire distribution.
\end{subequations}
